\paragraph{About this part.} Part~III listed what a faithful simulation of execution requires. Part~IV
shows one system built to meet those requirements, as a worked case study rather than an empirical
result. It describes the overall architecture (Section~\ref{sec:kaspar}); how simulated orders are
placed into, and filled from, a reconstructed queue of real orders (Section~\ref{sec:kaspar-fills});
how recorded exchange data are turned into replayable input (Section~\ref{sec:kaspar-pipeline}); why a
replay gives the same answer every time (Section~\ref{sec:kaspar-determinism}); how the same code runs
in replay, paper and live trading (Section~\ref{sec:kaspar-equivalence}); a simple execution algorithm
that serves as a model-free baseline (Section~\ref{sec:shadow}); and how latency is measured
(Section~\ref{sec:kaspar-measurement}).

% ================================================================
\section{Architecture overview}
\label{sec:kaspar}
% ================================================================

This part describes one system that implements the chain~\eqref{eq:chain} from market data to
simulated or live execution: Kaspar, a trading and simulation system for CME futures (ES, NQ and
treasury futures) built on a C++20 actor framework~\citep{maciejewski_actor}. It is included as a
case study of the methodology argued for in Part~III, not as an empirical contribution; every
statement below is a description of the design and code, and the few numbers quoted come from
previously published measurements, attributed where they appear.

The data path is
\begin{center}\small
CME MDP~3.0 multicast $\to$ socket reader $\to$ message buffer $\to$ decoder $\to$ feed handler
$\to$ order book $\to$ strategy $\to$ simulated order manager $\to$ [exchange gateway $\mid$ simulated fill].
\end{center}
Every component is an actor: it owns its state and communicates only by messages. Actors can be
placed in a \emph{group}, which runs several actors on one thread, processing one message at a time
in arrival order. The simulated order manager (SOM) sits between the strategy and the venue. In
live mode it forwards orders to the exchange gateway (CME iLink~3); in simulation mode it sends
them into the reconstructed book, where they are matched against recorded order flow.

\begin{figure}[ht]
\centering
\begin{tikzpicture}[
  node distance=4mm and 6mm,
  box/.style={draw, rounded corners=2pt, minimum height=7mm, align=center, font=\small},
  arr/.style={-{Stealth[length=2mm]}, thick}]
\node[box] (md) at (0,0.7) {recorded MBO\\(replay)};
\node[box] (live) at (0,-0.7) {live MDP 3.0\\multicast};
\node[box] (ob) at (3.2,0) {order book\\(per order)};
\node[box] (strat) at (5.6,0) {strategy\\(lights)};
\node[box] (som) at (7.6,0) {SOM};
\node[box] (sim) at (10.6,0.9) {fill in\\reconstructed queue};
\node[box] (ilink) at (10.6,-0.9) {iLink 3\\gateway};
\draw[arr] (md) -- (ob);
\draw[arr] (live) -- (ob);
\draw[arr] (ob) -- (strat);
\draw[arr] (strat) -- (som);
\draw[arr] (som) -- node[above left, font=\scriptsize]{sim} (sim);
\draw[arr] (som) -- node[below left, font=\scriptsize]{live} (ilink);
\draw[arr, dashed] (sim.north) |- ++(0,5mm) -| node[pos=0.25, above, font=\scriptsize]{simulated order enters the book} (ob.north);
\end{tikzpicture}
\caption{Kaspar data and order paths. Everything downstream of the market-data source is the same
code in replay and live operation; the branch between simulated and live execution is a single
decision in the simulated order manager.}
\label{fig:kaspar}
\end{figure}

% ================================================================
\section{Queue-exact simulated fills}
\label{sec:kaspar-fills}
% ================================================================

The central requirement of Section~\ref{sec:replay} is that a simulated order should be filled only if
a real order in its place would have been. This section explains how Kaspar meets it: simulated
orders sit in the same per-order queue as real ones, reach the book after a modelled delay, and are
filled only by real executions that reach them. Here and throughout, \emph{queue-exact} means exact
with respect to the displayed order sequence carried by the input MBO feed and the documented matching
rules; it does not imply knowledge of hidden liquidity or unobserved participants, and it does not model
the market impact of the simulated orders.

\subsection{The book holds every order}

The simulator reconstructs a full market-by-order book from recorded exchange data. Each price
level is a queue that holds every resting order individually, in arrival order, and simulated
orders are held in the same queue alongside real ones. There is no estimated ``volume ahead''
variable: a simulated order's queue position is its position in the list. This is the property
that makes a fill inference definable at all~\citep{maciejewski_shadow}.

\subsection{Placement and latency}

In simulation mode the SOM turns a new order into a book \emph{add} message tagged with a
simulated venue identifier and sends it to the order-book actor. The book does not insert it
immediately. It holds the order on a delay queue until its modelled arrival time at the matching
engine:
\begin{lstlisting}[language=C++]
auto order_engine_arrive_time =
    order_leave_time
      + uint64_t(feed_delay) * 1000
      + uint64_t(std::max(40, eff_delay)) * 1000; // 40 us floor
\end{lstlisting}
The outbound order delay, the outbound cancel delay and the inbound feed delay are independent
parameters. The feed delay is added because the strategy saw the market late by that amount, so a
round trip is feed plus order, not order alone. When the order arrives it is appended at the tail
of the queue at its price, behind all depth already resting there.

\subsection{The fill rule}

A recorded cancellation of a real order simply unlinks it, and every order behind it moves up. A
recorded execution triggers simulated fills only when the executing real order is the first real
order in the queue and a simulated order sits ahead of it:
\begin{lstlisting}[language=C++]
// fill sim orders only if the executing order is at the front
if (o == q->get_head_no_sim())
{
  fill_prev_sim_order();
  ...
}
\end{lstlisting}
Any remaining executed quantity is carried forward to earlier simulated orders. This is the
price--time-priority argument of Section~\ref{sec:replay} in code: an aggressor that reached a
real order which arrived after the simulated order must have passed through the simulated order.
Simulated orders are excluded from the best bid and offer that the strategy sees, so a strategy
cannot react to its own quote. Market impact is not modelled~\citep{maciejewski_shadow}.

\begin{example}[A fill, step by step]
\label{ex:trace}
Bid queue at 5000.00 over seven successive steps \#0--\#6, sizes in lots, real orders R1--R4 and
one simulated order S:
\begin{lstlisting}
#0  resting:            R1(10) R2(5) R3(20)
#1  S(3) arrives:       R1(10) R2(5) R3(20) S(3)      appended at tail
#2  R2 cancels:         R1(10) R3(20) S(3)            S moves up
#3  exec 10 vs R1:      R3(20) S(3)                   no sim order ahead of R1: no fill
#4  exec 20 vs R3:      S(3)                          S is behind R3: no fill
#5  R4(8) joins:        S(3) R4(8)
#6  exec 5 vs R4:       R4 is head real, S ahead  ->  S filled 3 @ 5000.00
\end{lstlisting}
At step \#6 the aggressor reached R4, which arrived after S, so it must have passed through S. A
``touch'' backtest would have filled S at step \#3, three recorded events earlier and before the
queue ahead of it was consumed.
\end{example}

\paragraph{Market impact: a limitation of replay, and a possible extension.} Because the recorded
market behaves exactly as it did, whatever the simulated orders do, replay answers ``would my order have
been filled against this market?'' but not ``how would the market have moved because of my order?''. A
simulated order never consumes liquidity that real traders would have reacted to. Market impact is
therefore outside what this architecture measures, and it is not implemented.

One way to add a reacting market, which we describe as possible future work rather than a feature, is a
second market-data source: a queue-reactive simulator (Section~\ref{sec:qr}) in place of the recorded feed.
The separation between the source of market data and the book, order manager and strategy actors means
the rest of the pipeline would not change. The generator produces queue counts rather than individual
orders, so its events would have to be turned into synthetic order-level messages---order identifiers and
first-in-first-out placement---for the existing book. The researcher's orders would then change queue
sizes, the simulator's rates would respond, and the price response to the strategy's own trading could be
measured.

What such an extension would estimate is limited, and should be stated with it.
\begin{itemize}[leftmargin=1.5em,itemsep=1pt]
\item It would capture \emph{mechanical} impact: queues consumed by the strategy's orders, refilling at
  state-dependent rates, and price moves when a queue empties. Neural extensions of the queue-reactive
  model reproduce the square-root law of impact on Bund futures~\citep{deepqr}, so this part can be
  realistic.
\item It would not capture \emph{informational} impact: other participants inferring that a large buyer is
  present and repricing. In the queue-reactive model this sits in the random resets of the book, which do
  not respond to the strategy, and for large orders it is a substantial part of real impact.
\item Its impact would be a property of the simulator, and would need checking against impact measured on
  real trading---post-fill mark-outs from paper or live fills, or published impact laws---before being
  relied on.
\end{itemize}
The resulting division of labour would be: queue-exact replay for fill realism on the real market, without
impact; a reactive simulator for mechanical impact and for comparing tactics in a market that responds;
and paper and live trading as the ground truth for both.

% ================================================================
\section{The replay pipeline}
\label{sec:kaspar-pipeline}
% ================================================================

A traditional backtester reads bars or trades from a file and applies a fill rule to a strategy's
orders. A queue-exact replay must instead reproduce the market data exactly as the trading system
would have received it, so that the same code path builds the same book. Kaspar supports three market-data
sources, each feeding the same downstream code:

\begin{center}\footnotesize
\begin{tabularx}{\linewidth}{>{\raggedright\arraybackslash}X>{\raggedright\arraybackslash}X>{\raggedright\arraybackslash}p{0.2\linewidth}}
\toprule
Source & Reader & Clock \\
\midrule
Live CME MDP~3.0 multicast & socket reader & real time \\
Packet capture (\texttt{.pcap}) of the same multicast groups & PCAP reader & capture timestamps \\
Decoded order-by-order record file & binary file reader, directly into the books & file timestamps \\
\bottomrule
\end{tabularx}
\end{center}

\begin{figure}[ht]
\centering
\begin{tikzpicture}[
  node distance=3mm,
  box/.style={draw, rounded corners=2pt, minimum width=78mm, minimum height=6mm, align=center, font=\small},
  arr/.style={-{Stealth[length=2mm]}, thick}]
\node[box] (a) {recorded packet capture of MDP~3.0 multicast};
\node[box, below=of a] (b) {PCAP reader: Ethernet $\to$ IP $\to$ UDP, capture and hardware timestamps};
\node[box, below=of b] (c) {message buffer and decoder (Simple Binary Encoding templates)};
\node[box, below=of c] (d) {feed handler: incremental book updates, order by order};
\node[box, below=of d] (e) {order book with real and simulated orders in each FIFO queue};
\node[box, below=of e] (f) {strategy and execution actors (lights)};
\node[box, below=of f] (g) {simulated order manager: delays, then insertion at the queue tail};
\foreach \x/\y in {a/b,b/c,c/d,d/e,e/f,f/g} \draw[arr] (\x) -- (\y);
\draw[arr] (g.east) -- ++(7mm,0) |- node[pos=0.25, right, font=\scriptsize, align=left]{simulated\\orders} (e.east);
\end{tikzpicture}
\caption{The packet-level replay path. Every stage below the reader is the same code that runs on
live data; only the reader and the final branch of the order manager differ.}
\label{fig:pcap}
\end{figure}

The capture reader parses each packet down to the exchange payload, optionally strips a
hardware-timestamp trailer added by capture equipment, and hands the payload to the same message
buffer the live socket reader uses, with both the software and hardware timestamps. It reads one
packet per message to itself, so replay is an ordinary actor loop rather than a blocking file read.
Two limits are stated in the project documentation: capture replay performs no packet recovery or
A/B feed arbitration, so a gap in the capture propagates to the book; and the decoded-record path
bypasses decoding entirely and is the fast option for repeated backtests. In the distributed
version, the capture and record readers are complete as libraries but not yet selectable from the
trading binary's command line; the separate simulator program drives replay from decoded record
files.

The order book then holds every order individually (Section~\ref{sec:kaspar-fills}), so a simulated
order's queue position is exact, and fills are decided by the recorded executions that follow.

% ================================================================
\section{Determinism}
\label{sec:kaspar-determinism}
% ================================================================

In the replay configuration the file reader, the order book, the timer, the SOM, the strategy
actors and the position manager all run in one actor group on one thread. The reader is added to
the group last, so all consumers exist before the first market-data message is delivered. Three
design choices then make a replay repeatable:
\begin{enumerate}[leftmargin=1.5em]
\item \textbf{One message sequence.} A single thread processes every message---market data, orders,
  acknowledgements, fills, timer events---in one total order.
\item \textbf{Market time.} The timer actor is driven by market-data timestamps, not the wall
  clock, so timeouts fire at the same point in the data on every run.
\item \textbf{Seeded randomness.} Any randomised strategy decision uses a per-actor seeded
  generator set in configuration.
\end{enumerate}
The property this implies, and the test that would check it, is that replaying the same recorded
file twice with the same configuration yields byte-identical fill records. We argue the property
from the design; we do not report the test here.

Determinism addresses a common weakness of trading research: the research implementation and the
production implementation are often different programs, and differences between them are
discovered in production.

% ================================================================
\section{Replay, paper and live equivalence}
\label{sec:kaspar-equivalence}
% ================================================================

Two choices are independent: where market data come from (recorded file or live multicast), and
where orders go (reconstructed book or exchange). Everything downstream of the message buffer---the
decoder, feed handler, order book, strategy and SOM---is the same code whichever source feeds it.
The branch between simulated and live execution is a single condition in the SOM's new-order
handler; both paths apply the same pre-trade reject checks, with one additional rule in live mode
for cancels of orders not yet acknowledged by the exchange. The combination ``live data, simulated
execution'' is paper trading: orders enter the live reconstructed book and are filled by the live
order flow under the same rule as in replay.

The live order-entry path implements CME's iLink~3 protocol: binary message encoding, HMAC request
signing, sequence management, and failover between primary and secondary gateways. The project
documentation states that both the market-data handler and the order-entry session have passed CME's
automated certification. Strategies can be written as in-process actors in C++ or, through an
in-process foreign-function interface, in Rust.

Two qualifications keep the claim precise. Switching between simulated and live execution is a single
flag in the SOM, but in the distributed build that flag is set in the program source rather than in a
configuration file, and the distributed configuration runs in simulation mode; the equivalence
demonstrated day to day is therefore between replay and paper trading, with the live path sharing all
code up to the SOM branch. And equivalence of code is not equivalence of outcomes: replay assumes no
market impact (Section~\ref{sec:kaspar-fills}) and models latency with configured delays, so a strategy
that is profitable in replay can still behave differently live.

% ================================================================
\section{A model-free execution benchmark}
\label{sec:shadow}
% ================================================================

\pov{} is the passive execution algorithm used in the case study~\citep{maciejewski_shadow}. It
makes no prediction about book evolution. When a third-party limit order is added within a
configured distance of the touch, the algorithm may place its own order at the same price, records
that participant's exchange order identifier, and cancels when that order leaves the book. The
placement decision is randomised at a configured rate; cancellation can be delayed by a grace
period measured in events or time. The design comment for the grace period states the trade-off
of Section~\ref{sec:adverse} directly: cancel too quickly and fills are given up; cancel too
slowly and the order bears the adverse selection.

\begin{table}[!htbp]
\centering\small
\caption{Representative \pov{} parameters (one of the configurations in the code repository).}
\label{tab:shadow-params}
\begin{tabular}{lll}
\toprule
Parameter & Value & Meaning \\
\midrule
\texttt{place\_rate\_bp} & 300 & place on 3\% of qualifying third-party adds \\
\texttt{max\_dist} / \texttt{max\_dist\_cancel} & 10 / 12 & ticks from touch to place / to force cancel \\
\texttt{delayed\_cancel\_events} & 5 & grace period after the shadowed order leaves \\
\texttt{nlights\_per\_side} & 10 & number of lights per side; each holds one order at one price \\
\texttt{rng\_seed} & 1 & seed for the placement decision \\
\bottomrule
\end{tabular}
\end{table}

Its role in this survey is as a null model for passive placement. It defines the bottom rung of a
hierarchy:
\begin{center}
model-free placement $\;\to\;$ predictive placement $\;\to\;$ full predictive execution policy.
\end{center}
A predictive model that claims value for passive execution should beat the model-free benchmark
\emph{in the same replay, under the same fill rule}, on slippage and on post-fill mark-outs. Any
improvement measured against a naive backtest instead is confounded by the fill-inference bias
of Section~\ref{sec:replay}.

% ================================================================
\section{Measurement}
\label{sec:kaspar-measurement}
% ================================================================

Latency claims need their endpoints. In the published measurement of Kaspar's market-data
handler, the start point is a software timestamp at the socket read and the end point is the order
book being published; time spent in the network card and kernel before the read is outside the
measurement, which is therefore socket-to-book, not wire-to-book. Measured on live CME multicast,
the median is 8.0\,$\mu$s, with book-message medians of 7.1, 7.3 and 7.1\,$\mu$s for ES, NQ and ZN
and 99th percentiles of 18.5, 13.6 and 57.0\,$\mu$s~\citep{maciejewski_mdlatency}. These numbers are
quoted as an example of stating endpoints, not as a result of this paper.

The simulator's slippage instrumentation measures execution cost for each fill against four
references: the mid at decision time, the interval volume-weighted average price, same-side
passive fills by other participants, and the touch. Comparing against the interval VWAP separates
market drift from adverse selection. Post-fill mark-outs at fixed times and after a fixed number
of book events provide the mark-out $R^{\text{fill}}_{h}$ of~\eqref{eq:markout} (the event-count version is
its analogue on an event clock) and the adverse-selection quantity of Section~\ref{sec:adverse}.

The order-flow signals of Part~II---OFI, queue imbalance, micro-price, fill-probability
models---are not part of the system described here. They are the candidate inputs that
Part~V proposes to search over, using this replay as the evaluator.
